\documentclass[10pt,twocolumn,a4paper]{article}
\usepackage[margin=0.75in]{geometry}
\usepackage{cite}
\usepackage{amsmath,amssymb,amsfonts}
\usepackage{graphicx}
\usepackage{textcomp}
\usepackage{xcolor}
\usepackage{booktabs}
\usepackage{hyperref}
\usepackage{titlesec}

% Title formatting
\titleformat{\section}{\large\bfseries}{\thesection.}{0.5em}{}
\titleformat{\subsection}{\normalsize\bfseries}{\thesubsection.}{0.5em}{}

\date{}

\begin{document}

\title{\Large\bfseries Cross-Geo Product Matching and Grouping Using\\Multimodal Embeddings and Graph Clustering}

\author{
    \textbf{Richard Čerňanský} \\
    \textit{NI-ADM Competition 2026} \\
    \textit{FIT CVUT}\\
    9th of May 2026, Prague, Czechia
}

\maketitle

\begin{abstract}
\noindent We address the problem of identifying identical fashion products across 13 European markets, where the same 
item appears with different titles, prices, and images. Our system combines multilingual sentence embeddings [5], 
CLIP image embeddings [6], and structured metadata through a two-phase pipeline. Phase 1 classifies predefined groups 
using LightGBM [7] on aggregated pairwise features (F1 = 0.9745). Phase 2 performs open-world product 
clustering via FAISS [9] approximate nearest neighbors, a trained pair-level edge scorer, contrastive fine-tuning, and Leiden [4] 
graph clustering, achieving pairwise F1 = 0.9005 on items\_phase1. Inspired by the fashionID architecture [2], we train a multimodal 
projection that fuses text, image, and price signals into a shared 128-dimensional metric space. 
!!An ablation study confirms that each modality contributes and that the multimodal projection improves pairwise F1 by +0.10 over the text-only baseline.
\end{abstract}

\vspace{1em}

\section{Introduction}
Product matching is a key capability for e-commerce platforms, enabling de-duplication, competitive intelligence, and 
recommendation of substitutes [2]. GLAMI, a fashion search engine operating across 13 European markets with over 20 million items, faces this challenge at scale: the same physical product appears under different titles (in different languages and scripts), different prices (in local currencies from EUR to HUF to BGN), and often different images. As Toth et al. [2] observe, fashion product matching is more similar to person re-identification than matching other product categories, because visual information is highly important and products are rarely identified through unique names beyond brand names.

The NI-ADM 2026 competition formalizes this as two tasks. Phase 1 is \textbf{binary classification}: given 15,000 groups of 5 items, 
predict whether each contains at least one duplicate pair (evaluated by F1). Phase 2 is \textbf{open-world entity resolution}: 
given $\sim$200k unlabeled items, construct disjoint product groups (evaluated by pairwise F1). Three challenges make this 
problem difficult:
(a) \textbf{multimodality} - text, images, price, and categorical metadata must be fused effectively;
(b) \textbf{cross-geo variation} - $\sim$40\% of product labels span 2+ markets with titles in completely different languages;
(c) \textbf{class imbalance} - !!Phase 1 uses an 80/20 negative/positive split, and in Phase 2 the vast majority of item pairs 
are non-matches, making brute-force comparison of $\sim$20 billion pairs infeasible.

\section{Related Work}
\textbf{Entity matching.} Christophides et al. [10] survey end-to-end entity resolution for big data, identifying blocking,
 matching, and clustering as the three core stages. Mudgal et al. [8] explore deep learning architectures for entity matching 
 and show that hybrid models combining attribute-level and structure-level signals outperform single-signal approaches. 
 Peeters et al. [1] demonstrate that LLMs achieve very high precision but at costs prohibitive for millions of items. 
 Li et al. [11] show that pre-trained language models fine-tuned as binary classifiers on concatenated text pairs achieve 
 state-of-the-art on textual entity matching benchmarks. Peeters and Bizer [12] extend this to contrastive learning, 
 demonstrating more efficient use of training data than cross-encoder approaches.

\textbf{Multimodal embeddings.} Our text encoder builds on Sentence-BERT [5], which adapts pre-trained transformers 
to produce fixed-size sentence embeddings via siamese networks. The multilingual variant supports 50+ languages, 
critical for cross-geo matching. For images, we use CLIP [6], which learns a shared text-image embedding space through 
contrastive pre-training on 400M pairs. Toth et al. (Zalando paper) [2] find that CLIP encoders surprisingly outperform DINOv2 [13] on visual product matching despite DINO's image-only self-supervised training, and that image information dominates text for fashion matching. For scalable retrieval, Johnson et al. [9] introduce FAISS with IVF indexing for billion-scale approximate nearest neighbor search.

\textbf{The fashionID architecture.} Most relevant to our work is Toth et al. [2], who propose fashionID: a 
late-fusion encoder that concatenates pooled CLIP image embeddings, CLIP text embeddings, and a 3-dimensional 
numeric vector (price, log-price, number of sizes), then projects through a single linear layer into a 192-dimensional
metric space. Training uses supervised contrastive loss [14] with large batch sizes (16k) and batch sampling that maximizes positive pairs. A key design decision is freezing the pre-trained encoders and training only the linear projection (0.6M parameters), enabling training in under 10 minutes on a single GPU. Their ablation shows that numeric features contribute $\sim$1 AUCPR point. However, their approach has \textbf{limitations} in our setting:
(a) they rely solely on embedding distance without leveraging structured metadata like department or color IDs;
(b) they use brand-based blocking for retrieval, which is unavailable in the GLAMI dataset;
(c) a single cosine threshold cannot distinguish same-category products that have high semantic similarity but are distinct items.
We address these through a two-stage pipeline where FAISS provides candidates and a LightGBM classifier incorporates 
structured features for final edge decisions.

\textbf{Graph clustering.} Traag et al. [4] introduce the Leiden algorithm, which improves on Louvain by guaranteeing 
internally connected communities. This property is important for product grouping, where a single weak transitive edge 
in Louvain can mistakenly merge distinct product clusters.

\section{Methodology}

\subsection{Data Preparation}
The dataset comprises 928k labeled training items and $\sim$200k unlabeled items per phase across 13 markets. 
We parse departmentIds to integer sets and colorTagIdsString to integer sets for Jaccard computation. Prices are 
normalized to EUR using fixed exchange rates (CZK: 0.04, HUF: 0.0025, BGN: 0.5112, RON: 0.2, PLN: 0.2326). 
EDA confirmed 91.9\% of product labels share identical departmentIds across items; brandEditionTagId was dropped 
(99.7\% missing). Following [2], we build a numeric feature vector from log-normalized EUR price.

\subsection{Feature Extraction}
\textbf{Text embeddings.} Each item's title and truncated description (500 chars) are 
encoded using \textit{paraphrase-multilingual-mpnet-base-v2} [5] (768-dim, 50+ languages), L2-normalized.

\textbf{Image embeddings.} Product images are encoded using CLIP ViT-B/32 [6] (512-dim), with 100\% coverage. 
Following [2], who find that CLIP provides a stronger visual baseline than DINOv2, we use CLIP as our image encoder throughout.

\textbf{Contrastive fine-tuning.} We fine-tuned the text encoder with MultipleNegativesRankingLoss [12] on $\sim$450k positive pairs (same-label items) for 3 epochs (batch size 64, warmup ratio 0.1), shifting embeddings from generic semantic similarity toward product identity.

\subsection{Phase 1: Group Classification}
For each group of 5 items, we compute all 10 pairwise comparisons, extracting per pair: text cosine similarity, 
image cosine similarity, department Jaccard, color tag Jaccard, EUR price ratio (min/max), and same-geo flag. 
These are aggregated into several group-level features (max, mean, min, 2nd-highest for embedding similarities; 
max, mean for structured features; plus interaction terms). A LightGBM [7] classifier (300 trees, depth 5, lr 0.05, 
scale\_pos\_weight=4.0 for the 80/20 class imbalance) is trained with 5-fold stratified cross-validation on 15,000 
synthetic validation groups.

\subsection{Phase 2: Product Clustering}
Phase 2 requires clustering $\sim$200k items into product groups, evaluated by pairwise F1. Our pipeline follows the 
two-stage approach of [2] (encoding + retrieval) but extends it with learned edge scoring and graph-based clustering:

\textbf{(1) Candidate generation.} FAISS IndexIVFFlat [9] (400 Voronoi cells, 40 probes) finds top-20 nearest neighbors per item by inner product, reducing the search space from $\sim$20 billion to $\sim$2M pairs.

\textbf{(2) Edge scoring.} A pair-level LightGBM (500 trees, depth 6) scores each candidate using 6 features: text\_sim, img\_sim, price\_sim, dept\_jaccard, color\_jaccard, same\_geo. This replaces the raw cosine threshold used by [2], allowing incorporation of structured metadata signals.

\textbf{(3) Leiden clustering.} A weighted NetworkX graph is built from scored edges. Leiden [4] community detection identifies product groups with guaranteed internal connectivity.

\textbf{(4) Post-processing.} Clusters exceeding 150 items are recursively re-clustered at higher resolution.

\subsection{Multimodal Projection}
Inspired by fashionID [2], we train a projection head that fuses all modalities into a 128-dim metric space. The architecture: text\_proj (768 $\rightarrow$ 128, ReLU), img\_proj (512 $\rightarrow$ 128, ReLU), numeric (log-normalized price, has\_image flag, placeholder), fusion (259 $\rightarrow$ 128), L2-normalized. Unlike [2] who use a single linear layer with 192-dim output, we use ReLU activations and a smaller output dimension. Training uses InfoNCE loss with temperature 0.07 and batch sampling that maximizes positive pairs per batch (256 groups, 4 items per group). We also experimented with joint training of the text encoder and projection head simultaneously, where the encoder learns representations that complement the frozen CLIP image signal. !!For confidence estimation, we compute mean edge weight (LightGBM match probability) within each cluster.

\section{Baselines}
We compare our method against baselines for both phases. \textbf{For Phase 1}:
\textbf{B1 (Embedding threshold):} Thresholding max pairwise cosine similarity of multilingual text 
embeddings at 0.75 (F1 = 0.900). This handles cross-lingual matching via the multilingual encoder but 
ignores price, image, and metadata.

\textbf{B2 (LightGBM, text + structured):} LightGBM on 13 text and structured features without 
image embeddings (F1 = 0.914). Adding CLIP image features yields F1 = 0.964 (our full Phase 1 model).

\textbf{For Phase 2}:
\textbf{B3 (Text cosine + threshold):} FAISS with raw cosine threshold on paraphrase-multilingual-mpnet 
embeddings, Leiden clustering (pairwise F1 $\sim$ 0.81). This is the straightforward application of [2]'s retrieval approach without any fine-tuning or structured features.

\section{Experiments}

\subsection{Validation Strategy}
For Phase 1, we constructed 15,000 synthetic validation groups from labeled training items (3,000 positive, 12,000 negative), 
matching the 20/80 ratio. 5-fold stratified CV provides reliable estimates. For Phase 2, we sample 200k items from 
training data and evaluate pairwise F1 against ground-truth labels. The competition evaluates Phase 2 on a blind test set (items\_phase\_2.csv).

\subsection{Hyperparameter Selection}
Phase 1 LightGBM: 300 trees, depth 5, lr 0.05, 31 leaves, scale\_pos\_weight=4.0. Classification threshold swept in 0.01 
steps to maximize F1. For Phase 2, grid search over k in \{10, 20, 30\}, edge threshold in \{0.80, 0.85, 0.90\}, and Leiden 
resolution in \{0.8, 1.0, 1.2\} selected k=20, resolution=1.0. The edge scorer LightGBM threshold was tuned on held-out 
pairs (best at 0.88). Contrastive fine-tuning used 3 epochs, batch size 64, warmup ratio 0.1. The projection head was 
trained with temperature 0.07, cosine annealing LR schedule, and 5 epochs.

\subsection{Quantitative Results}

\begin{table}[htb]
\caption{Results across both phases. Phase 1 metrics from 5-fold CV (rows 1-3) and leaderboard (row 4). Phase 2 metrics 
from the competition server. P/R unavailable for server-evaluated submissions.}
\vspace{0.5em}
\centering
\resizebox{\columnwidth}{!}{%
\begin{tabular}{llccc}
\toprule
\textbf{Phase} & \textbf{Method} & \textbf{F1} & \textbf{P} & \textbf{R} \\
\midrule
1 & B1: Text emb. threshold (0.75) & 0.900 & 0.910 & 0.889 \\
1 & B2: LightGBM (text + struct.) & 0.914 & 0.945 & 0.886 \\
1 & Full: LightGBM (text+img+struct.) & 0.964 & 0.989 & 0.941 \\
1 & Leaderboard (submitted) & 0.974 & - & - \\
\midrule
2 & B3: Text cosine + Leiden & 0.810 & - & - \\
2 & + Contrastive FT + edge scorer & 0.883 & - & - \\
2 & + Multimodal projection (full) & 0.901 & - & - \\
2 & Top of leaderboard at the time of reporting & 0.942 & - & - \\
\bottomrule
\end{tabular}
}
\end{table}

\section{Results}

\subsection{Leaderboard Performance}
Phase 1 leaderboard F1 = \textbf{0.974}. Phase 2 pairwise F1 = \textbf{0.9005}, with a gap of $\sim$0.04 to the top score of 
0.9424 in the time of reporting. The Phase 2 progression from baseline (0.81) to final submission (0.9005) demonstrates the 
cumulative effect of contrastive fine-tuning (+0.07), the LightGBM edge scorer (+0.01 on top of fine-tuning), and the multimodal 
projection (+0.02).

\subsection{Ablation Study}
\textbf{Phase 1 ablation.} Starting from text-only embedding threshold (F1 = 0.900): adding structured features (price, dept, color) 
with LightGBM raised F1 to 0.914 (+0.014); adding CLIP image embeddings raised F1 to 0.964 (+0.050). Feature importance: 
img\_max (1041 splits), text x image interaction (745), emb\_max (592), price\_max (566). Notably, dept\_max has near-zero 
importance in Phase 1 due to low variance (Jaccard mean = 0.997), but department remains valuable as a feature in the Phase 2 
edge scorer.

\textbf{Phase 2 ablation.} The baseline text encoder with cosine thresholding yields pairwise F1 = 0.81. Contrastive fine-tuning (MNR loss, 450k pairs) shifts the embedding space toward product identity, bringing F1 to $\sim$0.87. Adding the LightGBM edge scorer (6 features replacing raw cosine as edge weight) improves to 0.8828. The multimodal projection, which fuses text, image, and price into a 128-dim space using InfoNCE loss, brings the final score to 0.9005. Each component provides a meaningful gain, confirming the value of both domain-specific fine-tuning and multimodal fusion.

\subsection{Strengths and Weaknesses}
\textbf{Strengths.} The pipeline successfully handles cross-lingual matching (40\% of labels span 2+ geos). Currency 
normalization turns price into a strong discriminator. The modular architecture allows independent improvement of each stage. 
The two-stage scoring (FAISS + LightGBM) overcomes the limitation of [2]'s single-threshold approach by incorporating structured 
metadata.

\textbf{Weaknesses.} The gap to the leaderboard top (0.04 pairwise F1) suggests room for improvement. The two-stage training (encoder then projection) is suboptimal: the text encoder never sees the image signal during fine-tuning. Confidence calibration shows only $\sim$67\% purity for clusters with mean edge weight $> 0.95$. We use smaller pre-trained models (ViT-B/32) than [2]'s ViT-bigG-14, which limits visual feature quality.

\section{Conclusion}
We presented a multimodal product matching system that combines text, image, price, and categorical metadata for cross-geo fashion product identification. Key findings:
(1) image embeddings provide the largest single-component improvement in Phase 1 (+0.050 F1);
(2) contrastive fine-tuning is the largest single improvement in Phase 2 (+0.07 pairwise F1);
(3) the multimodal projection inspired by fashionID [2] yields a further +0.02 by fusing all signals;
(4) a LightGBM edge scorer incorporating structured features outperforms raw cosine thresholding by enabling lower candidate thresholds while maintaining precision;
(5) Leiden clustering outperforms Louvain by producing tighter, internally connected communities.

\textbf{Future work.} (1) \textbf{Joint encoder + projection training:} currently the pipeline trains in two stages (fine-tune encoder, then freeze and train projection). Joint end-to-end training with InfoNCE would let the encoder learn representations that complement the image signal directly, potentially gaining 1-3 pairwise F1 points. (2) Scaling to larger pre-trained encoders (ViT-L or ViT-bigG as in [2]) could improve visual feature quality. (3) Brand name extraction from titles as an exact-match feature. (4) !! Iterative cluster refinement via within-cluster pairwise verification.

\begin{thebibliography}{00}
\bibitem{1} R. Peeters et al., ``Entity Matching using Large Language Models,'' arXiv:2310.11244, 2023.
\bibitem{2} S. Toth et al., ``End-to-end multi-modal product matching in fashion e-commerce,'' arXiv:2403.11593, 2024.
\bibitem{3} V. Vancura et al., ``Scalable Linear Shallow Autoencoder for Collaborative Filtering,'' RecSys 2022.
\bibitem{4} V.A. Traag et al., ``From Louvain to Leiden: guaranteeing well-connected communities,'' Sci. Reports 9(1):5233, 2019.
\bibitem{5} N. Reimers and I. Gurevych, ``Sentence-BERT: Sentence Embeddings using Siamese BERT-Networks,'' EMNLP 2019.
\bibitem{6} A. Radford et al., ``Learning Transferable Visual Models From Natural Language Supervision,'' ICML 2021.
\bibitem{7} G. Ke et al., ``LightGBM: A Highly Efficient Gradient Boosting Decision Tree,'' NeurIPS 2017.
\bibitem{8} S. Mudgal et al., ``Deep Learning for Entity Matching: A Design Space Exploration,'' SIGMOD 2018.
\bibitem{9} J. Johnson et al., ``Billion-scale similarity search with GPUs,'' IEEE Trans. Big Data 7(3):535-547, 2021.
\bibitem{10} V. Christophides et al., ``An Overview of End-to-End Entity Resolution for Big Data,'' ACM Comput. Surv. 53(6), 2020.
\bibitem{11} Y. Li et al., ``Deep entity matching with pre-trained language models,'' PVLDB 14(1):50-60, 2020.
\bibitem{12} R. Peeters and C. Bizer, ``Supervised Contrastive Learning for Product Matching,'' WWW Companion 2022.
\bibitem{13} M. Oquab et al., ``DINOv2: Learning Robust Visual Features without Supervision,'' arXiv:2304.07193, 2023.
\bibitem{14} P. Khosla et al., ``Supervised Contrastive Learning,'' NeurIPS 2020.
\end{thebibliography}


\end{document}